\section{3. Visión y viabilidad del
producto}\label{visiuxf3n-y-viabilidad-del-producto}

\subsection{3.1 Problema y contexto del
cliente}\label{problema-y-contexto-del-cliente}

Servicio Automotriz Hermanos Ruiz, S.A. de C.V. es un taller que atiende
vehículos ligeros (automóviles particulares, motocicletas y camionetas)
de clientes particulares y de pequeñas flotillas empresariales. Tiene 4
bahías de trabajo y un estacionamiento para 5 vehículos en espera. Hoy
las citas se apuntan en un cuaderno, las órdenes se llenan en papel y la
comunicación con el cliente es por teléfono o WhatsApp (ver el proceso
actual en 3.5).

Según el anteproyecto y la entrevista con el cliente del 19/10/2026, los
problemas principales son:

\textbf{• Disputas por cobros:} las autorizaciones son verbales o por
mensaje y no queda registro de qué se aprobó ni por cuánto («dijo que
sí, pero cobra otra cosa»).

\textbf{• Reclamos por daños:} no hay checklist ni fotografías del
estado en que llega el vehículo.

\textbf{• Exceso de llamadas:} el cliente llama para saber el avance y
la recepcionista interrumpe su trabajo para averiguarlo.

\textbf{• Poca visibilidad:} la carga de trabajo y las bahías se asignan
de palabra, y hay vehículos detenidos sin justificación.

\textbf{• Historial disperso:} las órdenes se archivan en carpetas y es
difícil consultar reparaciones y garantías anteriores.

Hallazgos de la entrevista que no estaban en el anteproyecto:

• El cliente quiere una doble verificación con código de 6 dígitos al
autorizar cotizaciones, en forma remota o presencial (RN-02).

• El cliente puede aprobar o rechazar por separado los conceptos
opcionales de una cotización (RN-03).

• Si una refacción tarda varios días, el vehículo pasa al
estacionamiento para liberar la bahía (RN-06).

• Las citas pedidas en línea no apartan lugar hasta que la recepcionista
las confirma (RN-08).

• Varios mecánicos pueden colaborar en una orden, pero siempre hay un
responsable (RN-09).

Además, hay cortes de luz e internet, por lo que el personal necesita
seguir capturando sin conexión, y el jefe de mecánicos trabajará con 1 o
2 tabletas en un ambiente con ruido, polvo y manos sucias (4.3).

\subsection{3.2 Objetivos del producto e
indicadores}\label{objetivos-del-producto-e-indicadores}

{\def\LTcaptype{none} % do not increment counter
\begin{longtable}[]{@{}
  >{\raggedright\arraybackslash}p{(\linewidth - 8\tabcolsep) * \real{0.2349}}
  >{\raggedright\arraybackslash}p{(\linewidth - 8\tabcolsep) * \real{0.2030}}
  >{\raggedright\arraybackslash}p{(\linewidth - 8\tabcolsep) * \real{0.1710}}
  >{\raggedright\arraybackslash}p{(\linewidth - 8\tabcolsep) * \real{0.1708}}
  >{\raggedright\arraybackslash}p{(\linewidth - 8\tabcolsep) * \real{0.2202}}@{}}
\toprule\noalign{}
\begin{minipage}[b]{\linewidth}\raggedright
\textbf{Objetivo}
\end{minipage} & \begin{minipage}[b]{\linewidth}\raggedright
\textbf{Indicador}
\end{minipage} & \begin{minipage}[b]{\linewidth}\raggedright
\textbf{Situación actual}
\end{minipage} & \begin{minipage}[b]{\linewidth}\raggedright
\textbf{Meta}
\end{minipage} & \begin{minipage}[b]{\linewidth}\raggedright
\textbf{Cómo se medirá}
\end{minipage} \\
\midrule\noalign{}
\endhead
\bottomrule\noalign{}
\endlastfoot
Eliminar las disputas por cobros no autorizados & \% de órdenes cerradas
cuya cotización tiene autorización digital registrada (fecha, hora, IP y
monto) & Autorizaciones verbales o por mensaje, sin registro & 100 \% de
las órdenes cerradas & Reporte mensual del sistema: órdenes cerradas con
autorización confirmada / total de órdenes cerradas \\
Reducir las llamadas de seguimiento a recepción & Llamadas por día de
clientes que preguntan por el avance de su vehículo & Por medir:
bitácora de llamadas de una semana antes de liberar el portal del
cliente & Al menos 50 \% menos que la medición inicial & Bitácora de
llamadas de recepción, una semana antes y una semana después de la
liberación \\
Contar con evidencia del estado del vehículo al recibirlo & \% de
recepciones con checklist completo y las 4 fotografías obligatorias &
Órdenes en papel sin checklist ni fotografías & 100 \% de las
recepciones & Reporte mensual del sistema de recepciones completas \\
\end{longtable}
}

\subsection{3.3 Alcance}\label{alcance}

{\def\LTcaptype{none} % do not increment counter
\begin{longtable}[]{@{}
  >{\raggedright\arraybackslash}p{(\linewidth - 4\tabcolsep) * \real{0.3333}}
  >{\raggedright\arraybackslash}p{(\linewidth - 4\tabcolsep) * \real{0.3333}}
  >{\raggedright\arraybackslash}p{(\linewidth - 4\tabcolsep) * \real{0.3333}}@{}}
\toprule\noalign{}
\begin{minipage}[b]{\linewidth}\raggedright
\textbf{Dentro del alcance de la Fase 1}
\end{minipage} & \begin{minipage}[b]{\linewidth}\raggedright
\textbf{Dentro del alcance de la Fase 2}
\end{minipage} & \begin{minipage}[b]{\linewidth}\raggedright
\textbf{Fuera de alcance}
\end{minipage} \\
\midrule\noalign{}
\endhead
\bottomrule\noalign{}
\endlastfoot
• Acceso del personal por rol (HU-27)

• Citas en línea y agenda con confirmación manual (HU-03, HU-04, HU-06)

• Recepción con checklist y fotografías, también sin conexión (HU-01,
HU-25, HU-26)

• Asignación de bahía y mecánico, y diagnóstico (HU-07 a HU-09)

• Cotización desglosada y autorización del cliente con código 2FA
(HU-02, HU-12 a HU-16)

• Solicitud de refacciones y proveedores (HU-10, HU-17, HU-18)

• Aviso de fin de trabajos, cobro, entrega y crédito a flotillas (HU-20
a HU-23)

• Portal del cliente responsivo y aviso de privacidad (HU-28, HU-30) & •
Búsqueda rápida de clientes y vehículos (HU-05)

• Tablero simplificado de bahías para tableta (HU-11)

• Consulta de existencias de refacciones (HU-19)

• Consulta de garantías por el cliente (HU-24)

• Panel y reportes gerenciales (HU-29, RF-11)

• Envío del enlace por SMS (RF-02.3) & • Facturación electrónica (CFDI)
y contabilidad

• Pagos en línea con tarjeta

• Aplicación móvil nativa (iOS o Android)

• Vehículos pesados, tráileres y maquinaria (RN-07)

• Nómina y control de asistencia del personal \\
\end{longtable}
}

\subsection{3.4 Stakeholders}\label{stakeholders}

{\def\LTcaptype{none} % do not increment counter
\begin{longtable}[]{@{}
  >{\raggedright\arraybackslash}p{(\linewidth - 8\tabcolsep) * \real{0.1922}}
  >{\raggedright\arraybackslash}p{(\linewidth - 8\tabcolsep) * \real{0.2031}}
  >{\raggedright\arraybackslash}p{(\linewidth - 8\tabcolsep) * \real{0.2349}}
  >{\raggedright\arraybackslash}p{(\linewidth - 8\tabcolsep) * \real{0.1496}}
  >{\raggedright\arraybackslash}p{(\linewidth - 8\tabcolsep) * \real{0.2202}}@{}}
\toprule\noalign{}
\begin{minipage}[b]{\linewidth}\raggedright
\textbf{Stakeholder}
\end{minipage} & \begin{minipage}[b]{\linewidth}\raggedright
\textbf{Rol en el negocio}
\end{minipage} & \begin{minipage}[b]{\linewidth}\raggedright
\textbf{Interés en el sistema}
\end{minipage} & \begin{minipage}[b]{\linewidth}\raggedright
\textbf{Poder / interés}
\end{minipage} & \begin{minipage}[b]{\linewidth}\raggedright
\textbf{Estrategia de participación}
\end{minipage} \\
\midrule\noalign{}
\endhead
\bottomrule\noalign{}
\endlastfoot
Dueños del taller & Propietarios; deciden la inversión y aprueban el
proyecto. & Reducir disputas por cobros, cuidar la rentabilidad y la
reputación, y no pagar licencias. & Alto poder -- alto interés &
Gestionar de cerca mediante entrevistas y revisiones de avance en cada
entrega, validación de requisitos y aprobación formal del alcance. \\
Gerente & Supervisa la operación y autoriza entregas a crédito a
flotillas. & Consultar tableros de control, controlar saldos y autorizar
créditos con respaldo. & Alto poder -- alto interés & Gestionar de cerca
mediante sesiones de definición de reglas de crédito y de reportes, y
pruebas de aceptación de su flujo. \\
Recepcionista & Registra ingresos, confirma citas y atiende a clientes.
& Reducir llamadas de seguimiento y tener el registro de recepción y
fotografías en un solo lugar. & Bajo poder -- alto interés & Mantener
informada e involucrada, participando en las pruebas de usabilidad,
capacitación y retroalimentación continua. \\
Jefe de mecánicos & Centraliza diagnóstico, cotización y asignación de
bahías en la tableta. & Flujos cortos, botones grandes y operación sin
conexión. & Alto poder -- alto interés & Gestionar de cerca los
prototipos en el entrono real (tableta, pruebas con manos sucias) y
ajustes por iteración. \\
Mecánicos de bahía y de apoyo & Realizan reparaciones, inspecciones y
cubrimientos de turnos. & Recibir órdenes claras y enfocarse en el
trabajo físico & Bajo poder -- alto interés & Mantener informado con
capacitaciones breves y un canal para reportar dificultades. \\
Encargado de almacén & Controla requisiciones, proveedores, costos e
inventario de consumibles. & Directorio de proveedores, registro de
costos y fechas, y tickets de compra por orden & Bajo poder -- alto
interés & Mantener informado e involucrado mediante entrevistas sobre su
flujo de compras y pruebas del módulo M6. \\
Clientes particulares & Dueños de autos, motocicletas y vehículos. & Ver
el estado de su vehículo, autorizar presupuestos y recibir un cobro
transparente & Bajo poder -- alto interés & Mantener informado con las
pruebas de usabilidad con algunos clientes y alguna encuesta breve de
satisfacción al finalizar el servicio. \\
Clientes de flotillas empresariales & Empresas con vehículos que pueden
recibir entrega a crédito. & Condiciones de crédito claras y seguimiento
de varios vehículos & Poder medio -- interés medio & Mantener
satisfechos al definir con el gerente las reglas del crédito y validar
el flujo con un representante. \\
Equipo de desarrollo & Diseñan, construyen y prueban la plataforma. &
Entregar un sistema que cumpla los requisitos y la norma ISO/IEC 25010 &
Alto poder -- alto interés & Gestionar de cerca con reuniones
periódicas, control de versiones y revisiones internas de calidad. \\
Docente evaluador & Revisa el documento y el producto en curso. & Que el
proyecto cumpla los criterios y entregas del curso & Alto poder -- alto
interés & Mantener satisfecho entregando a tiempo, pidiendo asesoría en
los avances y pidiendo aclaraciones de dudas. \\
\end{longtable}
}

\subsection{3.5 Proceso actual (as-is)}\label{proceso-actual-as-is}

\figura{figuras/proceso-as-is.png}{height=0.8\textheight}{3.1}{Diagrama de actividades del proceso actual (as-is) de
atención del taller, con los problemas identificados en cada etapa.
Elaboración propia}{docs/03-vision/proceso-as-is.puml}

El diagrama de actividades muestra el proceso actual de atención del
taller antes de implementar la solución tecnológica. Se organiza en
cuatro carriles (cliente, recepcionista, mecánico y registro manual) y
recorre el flujo completo: solicitud de cita, recepción del vehículo,
diagnóstico, cotización, autorización, reparación, entrega y archivo del
historial. Los recuadros rojos señalan los problemas detectados en cada
etapa, como el riesgo de errores en el registro de citas, los datos
incompletos en las órdenes en papel, la información dispersa e ilegible,
las demoras en la cotización manual, el exceso de llamadas y un
historial difícil de consultar.

\subsection{3.6 Estudio de viabilidad}\label{estudio-de-viabilidad}

{\def\LTcaptype{none} % do not increment counter
\begin{longtable}[]{@{}
  >{\raggedright\arraybackslash}p{(\linewidth - 4\tabcolsep) * \real{0.1708}}
  >{\raggedright\arraybackslash}p{(\linewidth - 4\tabcolsep) * \real{0.5513}}
  >{\raggedright\arraybackslash}p{(\linewidth - 4\tabcolsep) * \real{0.2779}}@{}}
\toprule\noalign{}
\begin{minipage}[b]{\linewidth}\raggedright
\textbf{Tipo}
\end{minipage} & \begin{minipage}[b]{\linewidth}\raggedright
\textbf{Análisis y evidencia}
\end{minipage} & \begin{minipage}[b]{\linewidth}\raggedright
\textbf{Conclusión (viable / viable con condiciones / no viable)}
\end{minipage} \\
\midrule\noalign{}
\endhead
\bottomrule\noalign{}
\endlastfoot
Técnica & La pila propuesta (Python con FastAPI, PostgreSQL y una PWA
con Service Workers e IndexedDB) es de código abierto y está bien
documentada. El mayor riesgo técnico es la operación sin conexión con
sincronización posterior (RNF-08, RNF-09, R-02), que se validará con un
prototipo en el Sprint 1. & Viable con condiciones: validar la
sincronización sin conexión en el Sprint 1. \\
Económica & Costo de licencias de \$0: todas las herramientas son de
código abierto o de planes educativos (RNF-07). El taller usa las
computadoras que ya tiene en recepción y almacén, y 1 o 2 tabletas para
el jefe de mecánicos. Los planes gratuitos de correo y almacenamiento
tienen límites (R-05, R-06). & Viable con condiciones: el taller aporta
las tabletas y el sistema respeta los límites de los planes
gratuitos. \\
Operativa & El personal tiene experiencia tecnológica baja a media
(4.3.1). Se atiende con flujos de pocos pasos y botones grandes (RNF-02,
RNF-03), capacitación breve y pruebas de usabilidad con el personal. Los
dueños y el gerente apoyan el proyecto porque reduce disputas y llamadas
(3.4). & Viable con condiciones: capacitar y probar con el personal
antes de la Demo final. \\
Legal & Se tratan datos de identificación y contacto de los clientes
(nombre, teléfono y correo) y datos de sus vehículos, incluidas
fotografías; no se tratan datos sensibles. Las finalidades y medidas de
seguridad (control de acceso por rol, HTTPS, tokens temporales y aviso
de privacidad) están en 4.8, conforme a la Ley Federal de Protección de
Datos Personales en Posesión de los Particulares (DOF 20-03-2025). En
desarrollo y pruebas solo se usan datos ficticios (RNF-05). & Viable \\
De calendario & El semestre tiene 7 sprints (S0 a S6) hasta la Demo
final de la semana 16. El backlog tiene 30 historias; el release v0.1.0
incluye el flujo núcleo y las historias Should o Could pueden pasar a la
Fase 2 si la velocidad del equipo no alcanza. & Viable con condiciones:
priorizar las historias Must. \\
\end{longtable}
}

\textbf{Conclusión general:} el proyecto es viable con condiciones:
validar en el Sprint 1 la operación sin conexión, que el taller aporte
las tabletas, respetar los límites de los planes gratuitos y priorizar
las historias Must del flujo núcleo.
